%!TEX program = xelatex
\PassOptionsToPackage{table}{xcolor}
\documentclass[a4paper]{oblivoir}
%%%%%%%%%%%%%%%%%%%%%%%%%%%%%%%%%%%%%%%%%%%%%%%%%%%%%%%%%%%%%%%%%%%%
%%% 전처리부 (preamble)
\usepackage{indentfirst}
\usepackage{graphicx}
\usepackage{tikz}
\usepackage{tikzpagenodes}
\usepackage[left = 2cm, right = 3cm, top = 2.5cm, bottom = 2cm, a4paper]{geometry}
\usepackage{lipsum}
\usepackage{fontspec}
\usepackage{kotex}
\usepackage{color}
\usepackage{xcolor}
\definecolor{skyblue}{RGB}{150, 210, 250}
\definecolor{skyblue2}{RGB}{60, 120, 250}
\colorlet{Mycolor1}{green!10!orange}
\definecolor{Mycolor2}{HTML}{00F9DE}
\usepackage{geometry}
\setmainhangulfont[Path=FONT/]{KOPUBWORLD_DOTUM_PRO_MEDIUM.OTF}

%%%%%%%%%%%%%%%%%%%%%%%%%%%%%%%%%%%%%%%%%%%%%%%%%%%%%%%%%%%%%%%%%%%%
% 문서 시작 (Document)

\begin{document}

\begin{flushleft}
    {\setmainfont[Path=FONT/]{KOPUBWORLD_DOTUM_PRO_BOLD.OTF} {\textcolor{blue!40}{{\fontsize{24pt}{28pt}\selectfont{CONTENTS}}}}}
\end{flushleft}
\begin{tikzpicture}
    \filldraw[fill=white, draw=white] (0, 0) rectangle (16.5, 0.5);
\end{tikzpicture}
\begin{tikzpicture}
    \filldraw[fill=blue!40, draw=blue!40] (0, 0) rectangle (16.5, 1.25);
    \node[draw = blue!40] at (2.3, 0.625)  {\setmainhangulfont[Path=FONT/]{KOPUBWORLD_DOTUM_PRO_MEDIUM.OTF}{\textcolor{white}{{\huge 인투더 대학수학}}}};
\end{tikzpicture}
\begin{tikzpicture}
    \filldraw[fill=white, draw=white] (0, 0) rectangle (16.5, 1.25);
\end{tikzpicture}
\begin{tikzpicture}
    \node[draw = white] at (2.3, 0.625) {\setmainfont[Path=FONT/]{KOPUBWORLD_BATANG_PRO_BOLD.OTF}{\textcolor{blue!40}{{\huge Theme 01 - Discrete Mathematics (이산수학)}}}};
\end{tikzpicture}

\noindent
\begin{tikzpicture}
    \filldraw[fill=blue!40, draw=blue!40] (0, 1) rectangle (16.5, 0.9);
\end{tikzpicture}

\begin{itemize}
    \setmainfont[Path=FONT/]{KOPUBWORLD_DOTUM_PRO_BOLD.OTF} \fontsize{15pt}{18pt} 
    \item Zero : \cnm{수학}, \cnm{수학 I}, \cnm{수학 II}의 기본 용어
    \item Graph : 함수의 그래프와 관련된 개념의 시각화
    \item Basic : Zero와 Graph의 내용을 수식으로 엄밀화
    \item Algebra : 2015 개정 교육과정 수학 I (2022 개정 교육과정 대수-Algebra)
    \item Calculus : 2015 개정 교육과정 수학 II (2022 개정 교육과정 미적분 I-Calculus)
\end{itemize}

\begin{tikzpicture}
    \filldraw[fill=white, draw=white] (0, 0) rectangle (16.5, 0.1);
\end{tikzpicture}
\begin{tikzpicture}
    \node[draw = white] at (2.3, 0.625) {\setmainfont[Path=FONT/]{KOPUBWORLD_BATANG_PRO_BOLD.OTF}{\textcolor{blue!40}{{\huge Theme 02 - Linear Algebra (선형대수학)}}}};
\end{tikzpicture}

\noindent
\begin{tikzpicture}
    \filldraw[fill=blue!40, draw=blue!40] (0, 1) rectangle (16.5, 0.9);
\end{tikzpicture}

\begin{itemize}
    \setmainfont[Path=FONT/]{KOPUBWORLD_DOTUM_PRO_BOLD.OTF} \fontsize{15pt}{18pt} 
    \item Zero : \cnm{수학}, \cnm{수학 I}, \cnm{수학 II}의 기본 용어
    \item Graph : 함수의 그래프와 관련된 개념의 시각화
    \item Basic : Zero와 Graph의 내용을 수식으로 엄밀화
    \item Algebra : 2015 개정 교육과정 수학 I (2022 개정 교육과정 대수-Algebra)
    \item Calculus : 2015 개정 교육과정 수학 II (2022 개정 교육과정 미적분 I-Calculus)
\end{itemize}

\thispagestyle{empty}
% \newpage

\begin{tikzpicture}
    \filldraw[fill=white, draw=white] (0, 0) rectangle (16.5, 0.1);
\end{tikzpicture}

\begin{tikzpicture}
    \node[draw = white] at (2.3, 0.625) {\setmainfont[Path=FONT/]{KOPUBWORLD_BATANG_PRO_BOLD.OTF}{\textcolor{blue!40}{{\huge Theme 03 - Analysis (해석학)}}}};
\end{tikzpicture}

\noindent
\begin{tikzpicture}
    \filldraw[fill=blue!40, draw=blue!40] (0, 1) rectangle (16.5, 0.9);
\end{tikzpicture}

\begin{itemize}
    \setmainfont[Path=FONT/]{KOPUBWORLD_DOTUM_PRO_BOLD.OTF} \fontsize{15pt}{18pt} 
    \item Zero : \cnm{수학}, \cnm{수학 I}, \cnm{수학 II}의 기본 용어
    \item Graph : 함수의 그래프와 관련된 개념의 시각화
    \item Basic : Zero와 Graph의 내용을 수식으로 엄밀화
    \item Algebra : 2015 개정 교육과정 수학 I (2022 개정 교육과정 대수-Algebra)
    \item Calculus : 2015 개정 교육과정 수학 II (2022 개정 교육과정 미적분 I-Calculus)
\end{itemize}

\newpage
\begin{tikzpicture}
    \filldraw[fill=white, draw=white] (0, 0) rectangle (16.5, 0.2);
\end{tikzpicture}
\begin{tikzpicture}
    \node[draw = white] at (2.3, 0.625) {\setmainfont[Path=FONT/]{KOPUBWORLD_BATANG_PRO_BOLD.OTF}{\textcolor{blue!40}{{\huge Theme 04 - Probability, Statistics (확률과 통계)}}}};
\end{tikzpicture}

\noindent
\begin{tikzpicture}
    \filldraw[fill=blue!40, draw=blue!40] (0, 1) rectangle (16.5, 0.9);
\end{tikzpicture}

\begin{itemize}
    \setmainfont[Path=FONT/]{KOPUBWORLD_DOTUM_PRO_BOLD.OTF} \fontsize{15pt}{18pt} 
    \item Zero : \cnm{수학}, \cnm{수학 I}, \cnm{수학 II}의 기본 용어
    \item Graph : 함수의 그래프와 관련된 개념의 시각화
    \item Basic : Zero와 Graph의 내용을 수식으로 엄밀화
    \item Algebra : 2015 개정 교육과정 수학 I (2022 개정 교육과정 대수-Algebra)
    \item Calculus : 2015 개정 교육과정 수학 II (2022 개정 교육과정 미적분 I-Calculus)
\end{itemize}

\thispagestyle{empty}
\end{document}